% !TEX root = ../main.tex

\chapter{Discussion}\label{chapter:discussion}

This chapter discusses the experimental findings in context, their limitations, and implications for deploying uncertainty-aware traffic prediction systems.

\section{Interpretation of Results}

\subsection{MC Dropout Uncertainty Quality}

The primary finding of this thesis is that MC Dropout with $S = 30$ stochastic forward passes achieves Spearman~$\rho = 0.4820$ between predicted uncertainty and absolute prediction error for Trial~8 (100 test graphs, 3,163,500 node-level predictions). This moderate rank correlation means that nodes assigned higher $\sigma$ by MC Dropout consistently tend to incur larger absolute errors ($p \approx 0$, $n = 3{,}163{,}500$). A graph-level bootstrap ($B = 10{,}000$) yields a 95\% confidence interval of $[0.460, 0.469]$ for the mean per-graph $\rho$, confirming that this estimate is statistically precise. In practice, this enables triage of predictions: high-uncertainty outputs can be flagged for manual review or full MATSim re-simulation, while low-uncertainty outputs can be accepted.

\textbf{Choice of $S = 30$.} We use $S = 30$ MC forward passes, following the widely adopted range of $S = 10$--$50$ in the literature~\cite{gal2016dropout, lakshminarayanan2017simple}. Our computational budget for 100 test graphs (3,163,500 nodes $\times$ 30 passes $=$ 94.9M forward evaluations) required approximately 228 minutes on a T4 GPU.  To verify that $S = 30$ is sufficient, we conducted a formal convergence analysis by re-running MC Dropout inference on 10 test graphs (316,350~nodes) with $S = 50$ forward passes, retaining all individual pass outputs.  Subsampling at $S \in \{5, 10, 15, 20, 25, 30, 35, 40, 45, 50\}$ yields the convergence curves in Figure~\ref{fig:s_convergence}.  Both aggregate $\rho$ and mean per-graph $\rho$ flatten sharply after $S \approx 25$: moving from $S = 30$ ($\rho = 0.458$) to $S = 50$ ($\rho = 0.463$) improves ranking quality by only 1.0\%, well within the per-graph $\rho$ variability (std $= 0.027$).  Mean uncertainty $\bar{\sigma}$ follows the same pattern, stabilising at $\approx 2.0$~veh/h for $S \geq 25$.  The marginal return of additional passes is negligible, confirming that $S = 30$ sits firmly on the plateau and represents a well-justified operating point.

\begin{figure}[htbp]
\centering
\includegraphics[width=\textwidth]{figures/t8_s_convergence.pdf}
\caption[$S$-convergence of MC Dropout uncertainty quality (T8)]{$S$-convergence of MC Dropout uncertainty quality for Trial~8 (10 test graphs, 316,350~nodes, weight-remapped model).  \textbf{(a)}~Spearman~$\rho$ between uncertainty $\sigma$ and absolute error $|e|$: both aggregate and mean per-graph $\rho$ plateau by $S \approx 25$--$30$.  The shaded band shows $\pm 1$ std across graphs.  \textbf{(b)}~Mean uncertainty $\bar{\sigma}$ stabilises by $S \approx 25$.  The circled point marks $S = 30$.  Source: \texttt{s\_convergence\_results.json}.}
\label{fig:s_convergence}
\end{figure}

\textbf{Aggregate vs per-graph $\rho$.} The aggregate Spearman~$\rho = 0.4820$ (computed over all 3,163,500 nodes) is higher than the mean per-graph $\rho = 0.464$ (computed within each of the 100 graphs and then averaged). This discrepancy arises because the aggregate computation benefits from inter-graph variation: scenarios with systematically higher errors also tend to have higher mean $\sigma$, inflating the aggregate rank correlation. The per-graph $\rho$ is the more conservative and appropriate measure of within-scenario uncertainty ranking quality, since in deployment each scenario is evaluated independently.

This result aligns with the theoretical motivation of Gal and Ghahramani~\cite{gal2016dropout}, who showed that MC Dropout approximates Bayesian inference by sampling from an approximate posterior over model weights.

\subsection{Cross-Trial Robustness of UQ Findings}

A key question is whether the UQ findings for the best-performing trial (T8) are specific to that configuration or reflect a general property of the architecture. Section~\ref{sec:t7_crosscheck} presents a full replication on Trial~7 ($\rho = 0.4460$, dropout $= 0.3$), which uses the same architecture but different hyperparameters. The results confirm that the core conclusions hold across both trials: selective prediction reduces MAE monotonically with uncertainty filtering (41.2\% at 50\% retention for T8 vs 38.3\% for T7), raw MC Dropout intervals severely undercover at all nominal levels, and conformal prediction restores near-exact nominal coverage for both models. Trial~7 exhibits a larger empirical scaling factor ($k_{95} = 16.15$ vs $11.34$), consistent with its higher dropout rate producing more diffuse but less calibrated uncertainty estimates. This cross-trial consistency indicates that the UQ findings are not artefacts of a single model configuration.

\subsection{Ensemble Results and Their Limitations}

Experiments~A and~B, evaluated with corrected weight remapping (Section~\ref{sec:ensemble_caveat}), show that MC~Dropout ($\rho = 0.4908$) outperforms seed-based ensemble variance ($\rho = 0.4370$) by 12.3\% and multi-model ensemble disagreement ($\rho = 0.4333$) by 13.3\%.  Combined MC~Dropout + ensemble variance ($\rho = 0.4909$) provides negligible improvement over MC~Dropout alone, indicating that ensemble variance is largely redundant once MC~Dropout is available.  The Experiment~A MC~Dropout result ($\rho = 0.4908$) is consistent with the standalone T8 evaluation ($\rho = 0.4820$); the slight improvement arises from averaging uncertainty estimates across 5 seeded runs, which reduces estimation noise.

Multi-model ensemble disagreement captures meaningful but weaker uncertainty information: the 5 models span individual $R^2$ from 0.5117 (T2) to 0.5957 (T8), and their inter-model disagreement correlates with prediction error at $\rho = 0.4333$.  The weighted ensemble prediction ($R^2 = 0.5656$) does not outperform the best individual model (T8, $R^2 = 0.5957$) because averaging across models of uneven quality (individual $R^2$ ranging from 0.5117 to 0.5957) dilutes the strongest predictor. This suggests that ensemble accuracy gains would likely require either greater architectural diversity or stronger constituent models than were available in the 1,000-scenario subset.

Whether architecturally diverse ensembles (e.g., combining GAT, GraphSAGE, and GCN models) would further improve upon these results is untested.

\subsection{Predictive Performance in Context}

Trial~8 achieves test $R^2 = 0.5957$ (MAE $= 3.96$ veh/h, RMSE $= 7.12$ veh/h) on the 80/10/10 split. These metrics reflect the model's ability to predict traffic volume \emph{changes} ($\Delta v$) at the node level across 3,163,500 predictions spanning 100 scenarios.

Trial~1, which uses a \texttt{Linear(64$\to$1)} final layer and zero dropout, achieves $R^2 = 0.7860$ and MAE $= 2.97$ veh/h---the highest accuracy among all trials. This demonstrates that the PointNetTransfGAT backbone learns expressive representations of the road network and policy inputs. However, Trial~1 is \emph{structurally incompatible} with MC Dropout uncertainty quantification: its zero effective dropout rate means all stochastic forward passes produce identical outputs ($\sigma = 0$ everywhere), making the uncertainty estimation framework undefined. The exclusion of Trial~1 from UQ experiments is therefore not a limitation but a necessary methodological constraint: MC Dropout requires non-zero dropout by definition~\cite{gal2016dropout}.

A plausible contributing factor to Trial~1's accuracy advantage is architectural: the \texttt{GATConv(64$\to$1)} final layer used in Trials~2--8 introduces a neighbourhood aggregation step at the output, which may smooth per-node predictions toward local averages. This smoothing introduces an inductive bias that is well-suited to spatially correlated phenomena but may be mismatched to the highly variable, node-specific $\Delta v$ targets in this dataset. A \texttt{Linear(64$\to$1)} projection, by contrast, produces the final prediction as a function solely of the node's own 64-dimensional embedding, without any neighbourhood-level influence at the output stage. However, Trial~1 also differs from Trials~2--8 in batch size (32 vs 8) and learning rate ($10^{-3}$ vs $5 \times 10^{-4}$), and these confounds cannot be isolated without dedicated ablation experiments. The architectural hypothesis therefore remains speculative in the absence of controlled comparisons.

The gap between Trial~1 ($R^2 = 0.79$) and Trial~8 ($R^2 = 0.60$) indicates room for improvement: an architecture that combines the accuracy of a simpler output head with a separate uncertainty mechanism (e.g., evidential deep learning~\cite{amini2020deep}) could potentially recover Trial~1-level accuracy while still providing uncertainty estimates. A heteroscedastic extension was investigated as Trial~9 but achieved $R^2 = 0.02$ due to variance inflation under the NLL objective~\cite{seitzer2022pitfalls} (Section~\ref{sec:trial9_heteroscedastic}); this confirms that direct uncertainty prediction for this task requires either additional constraints on the variance head or an alternative formulation beyond standard heteroscedastic NLL.

\textbf{Context:} All results in this thesis are based on 1,000 of 10,000 available MATSim scenarios (10\% subset). The original Natterer~et~al.~\cite{natterer2025ml} work used the full 10,000-scenario dataset and achieved $R^2 = 0.91$ overall. Our results represent a preliminary investigation under data-limited conditions; the 10$\times$ difference in training data scale is the most likely explanation for the performance gap.

\section{Conformal Prediction vs MC Dropout Intervals}

Conformal prediction provides \emph{coverage-guaranteed} intervals: the 95\% conformal interval ($\pm$14.68 veh/h) guarantees that at least 95\% of true values fall within it, on average, under the exchangeability assumption. This guarantee holds regardless of the model's internal uncertainty representation.

In contrast, MC Dropout uncertainty $\sigma$ is \emph{not} a calibrated standard deviation---a well-documented property of approximate Bayesian methods (Section~\ref{sec:reliability_diagram}). T8 requires a scaling factor of $k_{95} = 11.34$ to achieve 95\% coverage using $\pm k \sigma$ intervals---versus the $k = 1.96$ expected for a calibrated Gaussian. Even after temperature scaling ($T = 2.70$), the effective $k_{95} \approx 5.30$ remains far from 1.96, confirming that the raw $\sigma$ from MC Dropout is most effective as a ranking signal for selective prediction and error detection, rather than as a probability distribution. Adaptive conformal prediction (Section~\ref{sec:conformal_conditional}) combines the advantages of both approaches: by normalising nonconformity scores with $\sigma$, it produces uncertainty-adaptive intervals that maintain conditional coverage across uncertainty deciles while preserving the marginal guarantee.

These three approaches form a natural hierarchy: MC Dropout provides adaptive, node-specific uncertainty scores useful for prediction ranking and selective filtering; temperature scaling improves the calibration of the Gaussian interpretation; and conformal prediction provides formal coverage guarantees. In practice, combining MC Dropout (for selective prediction) with adaptive conformal intervals (for coverage guarantees) represents the most complete approach.

Table~\ref{tab:uq_summary} summarises the key metrics, strengths, and limitations of each UQ method evaluated in this thesis.

\begin{table}[htbp]
\centering
\caption[Summary comparison of UQ methods evaluated in this thesis]{Summary comparison of UQ methods evaluated on the T8 model (100 test graphs, 3,163,500 nodes unless noted). All methods are post-hoc and require no model retraining.}
\label{tab:uq_summary}
\footnotesize
\begin{tabular}{p{2.2cm} p{2.8cm} p{2.6cm} p{3.0cm}}
\toprule
\textbf{Method} & \textbf{Key Metric} & \textbf{Strength} & \textbf{Limitation} \\
\midrule
MC Dropout ($S\!=\!30$) & $\rho = 0.4820$; CRPS/MAE $= 0.857$ & Node-specific ranking; selective prediction $-$41.2\% MAE & Not calibrated ($k_{95} = 11.34$); $S$-fold cost \\[4pt]
Temperature scaling ($T\!=\!2.70$) & ECE: $0.265 \to 0.048$; PIT KS: $0.245 \to 0.104$ & Single-parameter; no retraining & Gaussian assumption; residual miscalibration at tails \\[4pt]
Standard conformal & 95\% PICP $= 95.01\%$; $q = 14.68$ veh/h & Coverage guarantee; distribution-free & Fixed-width; poor conditional coverage \\[4pt]
Adaptive conformal & Cond.\ coverage: $[90.0\%, 96.2\%]$ & Node-specific widths; conditional coverage & Requires $\sigma$ from MC Dropout \\[4pt]
Ensemble ($5\times$) & MC avg $\rho = 0.4908$; Ens.\ var.\ $\rho = 0.4370$; multi-model $\rho = 0.4333$ & Complementary to single-model MC & Weaker ranking than MC Dropout; no accuracy gain over best single model \\
\bottomrule
\end{tabular}
\end{table}

\section{Practical Implications}

\subsection{Deployment Recommendations}

The results presented in this thesis translate into concrete deployment guidance for practitioners implementing uncertainty quantification in GNN-based traffic prediction systems. We structure these recommendations into six categories, each addressing a distinct aspect of the deployment workflow:

\textbf{1. Uncertainty estimation method: Use MC Dropout with $S = 30$ samples.} MC Dropout achieves the strongest uncertainty-error correlation found in this study ($\rho = 0.4820$ standalone, $\rho = 0.4908$ averaged over 5 seeded runs), substantially outperforming seed-based ensemble variance ($\rho = 0.4370$, $+12.3\%$ improvement) and multi-model ensemble disagreement ($\rho = 0.4333$, $+13.3\%$ improvement). For Trial~8 evaluated on 3.16M node-level predictions, this requires approximately 228 minutes on a T4 GPU---manageable within typical policy evaluation timelines. A dropout rate of 0.2 (T8) outperforms 0.3 (T5--T7) in this setting. The $S$-convergence analysis (Figure~\ref{fig:s_convergence}) confirms that $S = 30$ sits firmly on the plateau: moving from $S = 30$ to $S = 50$ improves $\rho$ by only 1.0\%, well within per-graph variability. If deployment constraints require faster inference, $S = 20$ represents a reasonable trade-off, sacrificing approximately 2--3\% in uncertainty quality for a 33\% reduction in computational cost.

\textbf{2. Selective prediction for resource allocation: Filter by uncertainty percentile.} The most immediate practical benefit of uncertainty quantification is selective prediction---accepting confident predictions directly while routing uncertain ones to full MATSim re-simulation. Filtering to the 50\% most confident predictions (bottom 50\% by MC Dropout $\sigma$) reduces MAE by 41.2\% (from 3.95 to $\approx$2.32~veh/h). Even a conservative threshold that discards only the top 10\% most uncertain nodes yields an 18.3\% MAE reduction. This creates a natural tiered workflow: \emph{accept} predictions below the 50th uncertainty percentile with high confidence (expected MAE $\approx$2.32~veh/h); \emph{flag} predictions between the 50th and 90th percentiles for sensitivity analysis or cross-checking; \emph{reject} predictions above the 90th percentile and trigger full MATSim re-simulation. In deployment, this triage allows planners to allocate expensive simulation resources only where genuine uncertainty warrants verification, substantially reducing the overall computational burden of large-scale policy screening studies.

\textbf{3. Formal coverage guarantees: Use adaptive conformal prediction.} When prediction intervals with guaranteed coverage are required---for example, to ensure that reported confidence bounds are statistically valid---conformal prediction provides the principled solution. Standard conformal prediction achieves near-exact marginal coverage (90.02\% and 95.01\% at the respective nominal levels, Section~\ref{sec:conformal_results}) with fixed-width intervals ($\pm$14.68~veh/h for 95\% coverage on the primary 50/50 split). However, this uniform width over-covers low-uncertainty nodes (98.6\% observed vs 90\% nominal for decile D1) and under-covers high-uncertainty nodes (62.9\% for D10), as predicted by the theoretical impossibility of distribution-free conditional coverage~\cite{barber2021limits}. Adaptive conformal prediction resolves this limitation: by normalising nonconformity scores with MC Dropout $\sigma$, it produces node-specific interval widths that narrow conditional coverage variation from [62.9\%, 98.6\%] to [90.0\%, 96.2\%] across uncertainty deciles (Section~\ref{sec:conformal_conditional}). The combination of MC Dropout (for adaptive interval widths) with conformal prediction (for the marginal coverage guarantee) thus addresses both the ranking and calibration requirements simultaneously. Importantly, conformal prediction requires a calibration set: in this study, a 50/50 calibration/evaluation split of the 100-graph test set (1,581,750 calibration nodes, 1,581,750 evaluation nodes) was sufficient; for deployment, practitioners should reserve 20--50\% of a held-out dataset for conformal calibration, with larger calibration sets yielding tighter quantile estimates.

\textbf{4. Probabilistic calibration: Apply temperature scaling as an intermediate step.} When approximate calibrated probabilities are needed but formal conformal guarantees are not required, temperature scaling provides a lightweight post-hoc improvement. Optimising a single scalar parameter $T$ on a held-out calibration set (20 graphs, 632,700 nodes in this study), temperature scaling reduces Expected Calibration Error by 82\% (from ECE $= 0.269$ to $0.048$, Section~\ref{sec:reliability_diagram}) and substantially improves the PIT diagnostic (Kolmogorov--Smirnov statistic drops from 0.245 to 0.104). The optimal $T = 2.70$ for Trial~8 means that raw MC Dropout $\sigma$ should be multiplied by 2.70 before interpreting it as a Gaussian standard deviation. However, residual miscalibration persists even after scaling (observed coverage of 83.3\% vs 95\% nominal at the highest confidence level), confirming that temperature scaling is useful for improving calibration but should not replace conformal prediction when formal coverage guarantees are required. The calibration/evaluation split used here (20/80) differs from the conformal split (50/50) by design: temperature scaling optimises a single parameter and benefits from a larger evaluation set for precise ECE estimation, whereas conformal prediction estimates a data-dependent quantile and benefits from a larger calibration set.

\textbf{5. Interpreting raw MC Dropout uncertainty: Use $\sigma$ for ranking, not as a standard deviation.} A critical finding is that raw MC Dropout $\sigma$ requires a scaling factor of $k_{95} = 11.34$ (versus the ideal 1.96 for a calibrated Gaussian) to achieve 95\% empirical coverage. This confirms that $\sigma$ is not a calibrated probability distribution in absolute terms---a well-established limitation of approximate Bayesian inference noted by Gal and Ghahramani~\cite{gal2016dropout} and systematically documented by Laves~et~al.~\cite{laves2020well}. However, this miscalibration does \emph{not} diminish the value of $\sigma$ as a ranking signal: the Spearman~$\rho = 0.4820$ correlation, the selective prediction gains, and the AUROC results (0.76 for top-10\% error detection) all rely on the \emph{ordering} of $\sigma$ values, not their magnitudes. In deployment, practitioners should use MC Dropout $\sigma$ to rank predictions by reliability (e.g., for selective filtering or error flagging) rather than interpreting it directly as a confidence interval width. For calibrated intervals, apply either temperature scaling (for approximate improvement) or conformal prediction (for formal guarantees).

\textbf{6. Ensemble methods: Multi-model ensembles provide complementary but weaker uncertainty.} The corrected ensemble experiments (Section~\ref{sec:ensemble_caveat}) show that MC~Dropout ($\rho = 0.4908$) outperforms seed-based ensemble variance ($\rho = 0.4370$) and multi-model disagreement ($\rho = 0.4333$) on architecturally identical models. Combined MC~Dropout + ensemble variance ($\rho = 0.4909$) yields negligible improvement, indicating that ensemble variance is largely redundant once MC~Dropout is available. However, multi-model disagreement may still provide value in production systems where multiple candidate models are maintained for robustness: disagreement across models signals scenarios where architecture-specific inductive biases diverge, a form of epistemic uncertainty not captured by single-model MC~Dropout. Importantly, the weighted ensemble prediction ($R^2 = 0.5656$) does not improve over the best individual model (T8, $R^2 = 0.5957$), suggesting that accuracy gains from ensembling require either architectural diversity (not tested here) or higher individual model quality than achieved with the 1,000-scenario subset. For deployment under computational constraints, single-model MC~Dropout is the recommended starting point; ensembles are a secondary consideration for systems where model diversity is already maintained for other reasons.

\subsection{Decision Support Integration}

Uncertainty estimates translate naturally into a tiered review workflow for urban planners. Predictions in the bottom 50\% of MC Dropout $\sigma$ can be accepted with confidence, with an expected MAE of $\approx$2.32~veh/h. Those in the 50--90\% range are candidates for sensitivity analysis or cross-checking against alternative models. Predictions in the top 10\% of $\sigma$ should be flagged for full MATSim simulation verification before informing a policy decision. This tiered approach lets planners allocate expensive full-simulation resources only where genuine uncertainty warrants it.

\begin{figure}[htbp]
\centering
\includegraphics[width=\textwidth]{figures/fig9_policy_explanation.pdf}
\caption[Uncertainty-guided policy decision framework (T8)]{Uncertainty-guided decision framework. The GNN surrogate (T8, $S=30$ MC Dropout) routes each prediction: \textbf{ACCEPT} (bottom 50\% $\sigma$, MAE $\approx$2.32~veh/h, $-$41.2\%); \textbf{FLAG} (50--90\%, manual review); \textbf{REJECT} (top 10\%, full MATSim re-simulation). Source: \texttt{mc\_dropout\_full\_100graphs\_mc30.npz}.}
\label{fig:policy_explanation}
\end{figure}

\section{Limitations}

\subsection{10\% Data Subset}

All experiments in this thesis use only 1,000 of 10,000 available MATSim scenarios. This limits the diversity of training examples and likely contributes to the gap between Trial~8's $R^2 = 0.5957$ and the higher performance that might be achievable with the full dataset. How MC Dropout uncertainty quality scales with more training data remains untested.

\subsection{Ensemble Experiment Scope}

The ensemble experiments evaluate only architecturally identical models (all five use the PointNetTransfGAT architecture with varying hyperparameters).  A PyG \texttt{GATConv} API version mismatch originally corrupted these results; the corrected evaluation (Section~\ref{sec:ensemble_caveat}) restores valid predictions and enables a reliable comparison.  However, the ensemble is limited to models from the same architecture family---true architectural diversity (e.g., combining GAT, GraphSAGE, and GCN models) might unlock further UQ benefits not observed here.

\subsection{Single Case Study}

Our experiments use a single road network (Paris). Generalisation to other cities, network topologies, or policy types requires additional validation. The feature correlation patterns (e.g., high-traffic roads being harder to predict) may not hold universally.

\subsection{Epistemic vs Aleatoric Uncertainty}\label{sec:limitations_epistemic}

MC Dropout captures only \emph{epistemic} (model) uncertainty---the component that arises from limited training data and could in principle be reduced by training on more scenarios. It does not decompose the total predictive uncertainty into epistemic and \emph{aleatoric} (data) components. Aleatoric uncertainty in this setting would reflect inherent randomness in traffic assignment outcomes for the same input (e.g., stochastic route choice in MATSim). Decomposing the two would require a heteroscedastic output layer (e.g., predicting both $\mu$ and $\log\sigma^2$ per node) or an evidential deep learning approach~\cite{amini2020deep}. Without this decomposition, we cannot determine whether the model's uncertainty is dominated by data noise or by genuine knowledge gaps---a distinction that would inform whether collecting more training data or improving model architecture is the more productive direction. A heteroscedastic extension was investigated in this work as Trial~9 (Section~\ref{sec:trial9_heteroscedastic}): training on the same 1,000-scenario 80/10/10 split converged to $R^2 = 0.02$ after 54 epochs, indicating that the NLL objective caused the model to inflate predicted variance rather than improve mean accuracy---a known failure mode under limited model capacity~\cite{seitzer2022pitfalls}. This negative result confirms that epistemic--aleatoric decomposition remains an open challenge for this architecture; evidential deep learning~\cite{amini2020deep} or a regularised heteroscedastic formulation represent viable directions for future work.

\subsection{Absence of Non-GNN Baselines}

This thesis evaluates UQ methods exclusively on the PointNetTransfGAT architecture. No non-GNN baselines (e.g., random forests with quantile regression, Gaussian processes on graph embeddings, or simple MLP ensembles) are included. Such baselines would contextualise whether the observed $\rho = 0.4820$ reflects a strength of the GNN architecture or is achievable by simpler models. The omission reflects the scope of this work, which builds on an existing GNN pipeline~\cite{natterer2025ml} and focuses on post-hoc UQ for that specific model rather than a comparative model selection study.

\subsection{Calibration and Conformal Data Splits}

Temperature scaling and conformal prediction use different calibration/evaluation splits of the 100 test graphs. Temperature scaling uses the first 20 graphs for calibration and the remaining 80 for evaluation, while conformal prediction uses a 50/50 random split. This design choice was made to maximise the calibration set size for conformal prediction (which benefits from more calibration data for tighter quantile estimates) while keeping the temperature scaling calibration set smaller (since it optimises a single scalar parameter). However, this means the ECE and PICP metrics are not evaluated on identical held-out sets, and results should be compared with this caveat in mind.

\subsection{Computational Cost of Post-Hoc Methods}

Temperature scaling and conformal prediction are computationally inexpensive: both operate on pre-computed predictions and uncertainties and require no model re-inference. Temperature scaling involves a grid search over $\sim$100 candidate $T$ values, completing in under 10 seconds on the 632,700-node calibration set. Conformal prediction requires sorting 1,581,750 nonconformity scores and computing a single quantile, also completing in seconds. The dominant cost is the MC Dropout inference itself (228 minutes for 100 graphs on a T4 GPU), which is required by all UQ methods in this thesis. Post-hoc calibration adds negligible overhead.

\subsection{MC Dropout Calibration: A Known Limitation with Practical Solutions}

The $k_{95} = 11.34$ finding confirms that MC Dropout $\sigma$ is not a calibrated standard deviation---a property that is well established in the approximate Bayesian inference literature. Gal and Ghahramani~\cite{gal2016dropout} note that the variational posterior produced by MC Dropout is typically too concentrated, and subsequent work has consistently observed that raw MC Dropout uncertainty underestimates true prediction variance in absolute terms~\cite{laves2020well}. The reliability diagram analysis (Section~\ref{sec:reliability_diagram}) quantifies this systematically for our setting: the Expected Calibration Error (ECE) of 0.265 means that, on average, observed coverage deviates from the nominal level by 26.5 percentage points.

This miscalibration does \emph{not} diminish the ranking value of MC Dropout uncertainty. The Spearman~$\rho = 0.4820$ correlation, the selective prediction gains ($-$41.2\% MAE at 50\% retention), and the AUROC results (0.76 for top-10\% error detection) all rely on the \emph{ordering} of $\sigma$ values, not their absolute magnitudes. Post-hoc temperature scaling~\cite{guo2017calibration, laves2020well} with an optimal $T = 2.70$ reduces the ECE from 0.269 to 0.048---an 82\% improvement---confirming that a simple rescaling captures most of the calibration gap. However, residual miscalibration persists at high nominal levels (83.3\% observed vs 95\% nominal), reflecting the non-Gaussian shape of the residual distribution, which has heavier tails than any single-parameter Gaussian rescaling can capture. Temperature scaling is therefore a useful intermediate step between raw MC Dropout and full conformal prediction, but should not be relied upon for formal coverage guarantees. The three-method hierarchy demonstrated in this thesis---MC Dropout for ranking, temperature scaling for improved calibration, conformal prediction for formal coverage guarantees---provides a practical resolution that addresses each level of the calibration problem.

\textbf{Proper scoring rules and the PIT diagnostic.} Three additional metrics confirm and extend the calibration picture. The CRPS of 3.38~veh/h is 14.3\% below the MAE (CRPS/MAE $= 0.857$), indicating that the Gaussian predictive distribution provides genuine information beyond the point estimate even though the raw $\sigma$ is not calibrated in absolute terms. To contextualise this ratio, we note that for a perfectly calibrated Gaussian---where observations are truly drawn from $\mathcal{N}(\mu, \sigma^2)$---the expected CRPS equals $\sigma/\sqrt{\pi}$ while the expected absolute error equals $\sigma\sqrt{2/\pi}$, yielding a theoretical optimum of $\text{CRPS}/\text{MAE} = 1/\sqrt{2} \approx 0.707$ regardless of $\sigma$~\cite{gneiting2007strictly}. Our observed ratio of 0.857 exceeds this optimum by approximately 21\%, directly quantifying the calibration cost of the underdispersed MC~Dropout $\sigma$. In probabilistic weather forecasting---where CRPS is the standard evaluation metric---well-calibrated post-processed ensemble models achieve ratios near this theoretical optimum~\cite{gneiting2005calibrated}. Direct comparison to GNN uncertainty quantification literature is limited because most graph-based UQ work targets node classification rather than regression~\cite{fuchsgruber2024energy}, though Murad~et~al.~\cite{murad2021probabilistic} evaluate MC~Dropout on graph convolutional networks for spatial air quality prediction and similarly observe that CRPS closely tracks MAE for uncalibrated probabilistic models. The PIT histogram (Section~\ref{sec:pit_results}) makes the miscalibration visually explicit: a prominent spike near zero (28.4\% of observations in the first 5\% bin) and U-shaped tails reveal systematic underdispersion, consistent with the ECE of 0.265 and the scaling factor $k_{95} = 11.34$. Importantly, the PIT also exposes an \emph{asymmetry} not visible in the ECE: the left-skewed PIT mean (0.433 vs the ideal 0.500) indicates that the model overpredicts more often than it underpredicts at extreme confidence levels. Finally, the Winkler interval scores (Section~\ref{sec:winkler_results}) quantify the calibration penalty in decision-relevant terms: raw Gaussian 90\% intervals achieve only 48.6\% coverage and incur a mean Winkler score of 49.68, while conformal sigma-scaled intervals achieve the target 90\% coverage with a 35\% lower score (32.33). The sigma-scaled advantage over fixed-width conformal intervals (10\% improvement at 90\%) confirms that the MC~Dropout ranking signal retains value even within the conformal framework.

\textbf{Effect of temperature scaling on the PIT diagnostic.} Since the PIT histogram reveals severe underdispersion under raw $\sigma$, a natural question is whether temperature scaling alleviates it. Replacing $\sigma$ with $\sigma \cdot T$ ($T = 2.70$) in the PIT computation yields a substantially improved histogram (Figure~\ref{fig:pit_after_tempscaling}): the Kolmogorov--Smirnov statistic drops from 0.245 to 0.104 ($-$57\%), the PIT mean shifts from 0.433 to 0.471 (closer to the ideal 0.500), and the PIT standard deviation decreases from 0.399 to 0.302 (ideal: 0.289). The first-bin spike shrinks from 28.4\% to 8.8\%, eliminating the most extreme underdispersion artefact. However, the last bin remains elevated (11.7\% vs the ideal 5.0\%), indicating that the Gaussian assumption still cannot fully capture the heavy right tail of the residual distribution. This partial improvement is consistent with the residual ECE of 0.048 after scaling and reinforces the conclusion that temperature scaling is a useful intermediate step but not a substitute for distribution-free conformal prediction when formal coverage guarantees are required.

\begin{figure}[htbp]
\centering
\includegraphics[width=\textwidth]{figures/t8_pit_after_tempscaling.pdf}
\caption[PIT histograms before and after temperature scaling (T8)]{PIT histograms for T8 MC Dropout ($S = 30$, 3,163,500 nodes) before and after temperature scaling ($T = 2.70$). \textbf{Left:} Raw $\sigma$ produces severe underdispersion (KS $= 0.245$, first-bin spike at 28.4\%). \textbf{Right:} Scaling $\sigma$ by $T = 2.70$ reduces the KS statistic to 0.104 ($-$57\%) and brings the PIT distribution substantially closer to uniformity, though residual deviation in the last bin (11.7\%) indicates the Gaussian assumption remains imperfect. Source: \texttt{pit\_after\_tempscaling\_t8.json}.}
\label{fig:pit_after_tempscaling}
\end{figure}

\subsection{Conditional Coverage: A Theoretical Limitation Resolved by Adaptive Conformal Prediction}

Barber~et~al.~\cite{barber2021limits} prove that distribution-free conditional coverage is impossible without additional assumptions---a fundamental theoretical result that applies to \emph{all} conformal methods, not a weakness of our particular implementation. Our analysis confirms this theoretical prediction in practice: standard conformal intervals with a global quantile over-cover low-uncertainty nodes (98.6\% at 90\% nominal for decile D1) and under-cover high-uncertainty nodes (62.9\% for D10). This disparity is the mathematically inevitable consequence of applying fixed-width intervals to a heterogeneous prediction task.

Adaptive conformal prediction, which normalises nonconformity scores by $\sigma$, provides the principled solution: it produces node-specific interval widths that are wider where MC Dropout signals higher uncertainty and narrower where it signals confidence. This largely resolves the conditional coverage disparity (90.0\%--96.2\% across deciles at the 90\% level), demonstrating that the MC Dropout ranking signal carries genuine calibration information even when its absolute magnitudes are miscalibrated. The combination of MC Dropout uncertainty (for adaptive interval widths) with conformal prediction (for the marginal coverage guarantee) thus addresses both the ranking and calibration limitations identified in this thesis.

The stratified analysis by feature quartiles (Section~\ref{sec:stratified_uq}) reveals that prediction difficulty and uncertainty quality both vary substantially across the input space. Low-volume, low-capacity road segments are predicted with high accuracy (MAE $\approx 1.0$--$1.4$ veh/h) and excellent ranking ($\rho \approx 0.74$--$0.81$), while high-traffic segments are harder to predict (MAE $\approx 4.5$--$8.0$ veh/h) with degraded ranking ($\rho \approx 0.26$--$0.41$). This heterogeneity has practical implications for deployment: uncertainty-guided filtering is most effective for the easy-to-predict majority of the network, where the model already achieves high accuracy. For the most challenging high-volume segments, the uncertainty signal is weaker---motivating future work on feature-conditional or group-conditional conformal methods that could provide tighter guarantees for specific road categories.

\subsection{Stability Across Scenarios and Trials}

The per-graph analysis (Section~\ref{sec:per_graph_variation}) provides reassurance that the aggregate UQ quality metrics are not driven by outlier scenarios. Spearman~$\rho$ is consistent across 100 test graphs (mean = 0.464, std = 0.023, range [0.41, 0.51]), confirming that uncertainty ranking is a robust property of the model across diverse traffic scenarios. Cross-trial validation on T7 (Section~\ref{sec:t7_crosscheck}, Section~\ref{sec:t7_auroc}) further strengthens this conclusion: T7 achieves AUROC values of 0.742 and 0.715 for top-10\% and top-20\% error detection---slightly below T8 but following the same pattern.

\section{Future Directions}

\subsection{Conformal Risk Control}

Beyond fixed-coverage conformal prediction, conformal risk control~\cite{angelopoulos2023conformal} enables coverage guarantees on arbitrary loss functions. Combined with the adaptive conformal approach validated in Section~\ref{sec:conformal_conditional}, this could provide even tighter, uncertainty-adaptive intervals for traffic prediction.

\subsection{Advanced Calibration Methods}

Temperature scaling ($T = 2.70$) reduced ECE by 82\% but left residual miscalibration at high nominal levels due to the non-Gaussian residual distribution. More flexible calibration approaches---such as isotonic regression, Platt scaling with learned variance, or quantile recalibration~\cite{kuleshov2018accurate}---could potentially close this remaining gap without requiring model retraining.

\subsection{Architecturally Diverse Ensembles}

The corrected ensemble experiments (Section~\ref{sec:ensemble_caveat}) show that multi-model disagreement across five PointNetTransfGAT variants provides useful but weaker uncertainty ranking ($\rho = 0.4333$) than single-model MC~Dropout ($\rho = 0.4908$).  A natural extension is to explore ensembles combining architecturally diverse models (e.g., GAT, GraphSAGE, GCN) trained on the same data, as true architectural diversity may unlock ensemble UQ benefits not observed with identical architectures.

\subsection{Post-Hoc Epistemic Uncertainty via Energy-Based Models}

An alternative to MC Dropout for epistemic uncertainty estimation is the Graph Energy-Based Model (GEBM) approach proposed by Fuchsgruber~et~al.~\cite{fuchsgruber2024energy}. GEBM fits a lightweight energy-based model to the GNN's learned node representations and quantifies epistemic uncertainty through the energy score of test nodes relative to the training distribution---requiring only a single deterministic forward pass at inference time. This eliminates the $S$-fold computational overhead of MC Dropout (228~minutes for $S = 30$ in our setting) and avoids the variance-convergence trade-off studied in Figure~\ref{fig:s_convergence}. GEBM also explicitly targets epistemic uncertainty, addressing the limitation that MC Dropout does not decompose total uncertainty into epistemic and aleatoric components (Section~\ref{sec:limitations_epistemic}). Applying GEBM to the PointNetTransfGAT architecture as a post-hoc method---without retraining---would provide a direct comparison against MC Dropout on the same model and dataset, and could clarify whether the moderate Spearman~$\rho$ observed here reflects a fundamental limitation of the architecture or of the uncertainty method.

\subsection{Full 10,000-Scenario Training}

The most direct improvement would be training on the full 10,000-scenario dataset to assess whether model accuracy and uncertainty quality improve proportionally with data scale.

\subsection{Out-of-Distribution Detection}

Explicit OOD detectors could flag scenarios unlike any training example, complementing uncertainty estimation with categorical reliability assessment for novel policy interventions.

\subsection{Active Learning}

Uncertainty estimates could guide active learning: selecting high-uncertainty scenarios for additional MATSim simulation to expand training data where most needed.
